\chapter{Commuting Observables, Energy Phases, and Spin Precession}

\lectureoverview{We begin with the qualitative meaning of uncertainty. Two observables can be definite together only when they possess a complete common eigenbasis, and this is equivalent to their operators commuting. We then review unitary evolution and the Schr\"odinger equation, compare commutators with classical Poisson brackets, solve the time-dependent equation in the energy basis, and apply the result to a spin in a magnetic field. The final calculation gives precession of the spin expectation values, the two energy levels of a general spin Hamiltonian, and the distinction between a sideways expectation value and a photon of half the energy gap.}

The two-state spin remains our simplest quantum system. A three-state Hilbert space would be the next simplest abstract possibility; two spins already have four orthogonal basis states. We will not yet compose two spins, although that step is close and will bring entanglement with it. First we should finish the one-spin problem and put the ideas of uncertainty and time evolution on firmer ground.

\section{When Two Quantities Can Be Definite Together}

A good measurement is repeatable. If an observable is measured and the value \(l\) is obtained, an immediate repetition should return the same value. The first measurement must therefore leave the system, at least momentarily, in an eigenstate of the corresponding operator:
\begin{equation}
L\ket{l}=l\ket{l}.
\end{equation}
This operational statement is enough to expose the qualitative origin of uncertainty.

Suppose two observables \(L\) and \(M\) can be measured simultaneously. A joint measurement must leave both reported values definite. Its possible post-measurement states must consequently be simultaneous eigenvectors of both operators. In the finite-dimensional setting used here, we may choose a complete orthonormal set \(\{\ket{i}\}\) satisfying
\begin{equation}
L\ket{i}=l_i\ket{i},
\qquad
M\ket{i}=m_i\ket{i}.
\label{eq:l5-common-eigenvectors}
\end{equation}

\begin{remark}
The word \emph{basis} will mean a complete orthonormal basis. Linear algebra permits nonorthogonal bases, but the orthonormal convention is especially natural for quantum states. Degenerate eigenvalues leave freedom inside an eigenspace; within that space the eigenvectors can be chosen orthonormally.
\end{remark}

\begin{theorem}
Two finite-dimensional Hermitian observables are simultaneously measurable precisely when they admit a complete common eigenbasis, equivalently when they commute.
\end{theorem}

The forward direction is immediate from Equation~\eqref{eq:l5-common-eigenvectors}. Acting in one order gives
\begin{align}
ML\ket{i}
&=M\bigl(l_i\ket{i}\bigr)\\
&=l_i m_i\ket{i},
\end{align}
while the reverse order gives
\begin{align}
LM\ket{i}
&=L\bigl(m_i\ket{i}\bigr)\\
&=m_i l_i\ket{i}.
\end{align}
The eigenvalues are ordinary numbers, so \(l_i m_i=m_i l_i\). Thus \(ML\) and \(LM\) act identically on every member of a basis and hence on every state:
\begin{equation}
[L,M]\equiv LM-ML=0.
\label{eq:l5-commuting-condition}
\end{equation}

\begin{figure}[htbp]
\centering
\includegraphics[width=0.86\textwidth]{lecture_05_figure_02.png}
\caption{The common-eigenvector argument reaches \(LM-ML=[L,M]=0\).}
\lectureframe{00:10:30}
\end{figure}

The converse is also true: commuting Hermitian operators can be simultaneously diagonalized. If one operator has a degenerate eigenspace, the other preserves that eigenspace, and it can be diagonalized within it. Repeating this in every eigenspace produces a complete common eigenbasis.

The Pauli operators give the first important counterexample. No two of \(\sigma_x,\sigma_y,\sigma_z\) commute. The eigenstates of \(\sigma_z\) are up and down, while those of \(\sigma_x\) are right and left. One may prepare either component sharply and repeat its measurement, but one cannot prepare both components sharply at once. This is uncertainty before any numerical inequality has been introduced.

Later we will quantify how the uncertainties of two noncommuting observables constrain one another. No numerical uncertainty principle is being proved here. The present point is prior to the inequality: noncommutation prevents a state from assigning simultaneous definite values to both quantities.

\begin{classroomqa}
\audiencequestion{Why does a simultaneous measurement require common eigenvectors?}
\lecturerresponse{A repeatable measurement leaves a state with the measured value definite, which means an eigenstate of the measured operator. If two values are obtained together and both remain repeatable, the resulting state must be an eigenstate of both operators. The complete set of joint outcomes therefore supplies a common eigenbasis.}
\end{classroomqa}

\begin{classroomqa}
\audiencequestion{Does this involve two different particles?}
\lecturerresponse{It need not, as \(\sigma_x\) and \(\sigma_z\) already act on one spin. For two distinct subsystems, however, observables acting on different factors do commute: \(A\otimes I\) commutes with \(I\otimes B\). That is why a component of one spin and a component of another can be measured together. Composing such systems will later lead to entanglement.}
\end{classroomqa}

\section{Unitary Evolution and Its Generator}

Quantum mechanics took considerably longer than classical mechanics to define a state, but once the state space is established the rule for motion is compact. Let
\begin{equation}
\ket{\psi(t)}=U(t)\ket{\psi(0)}.
\label{eq:l5-propagator}
\end{equation}
Time evolution must preserve distinguishability. More strongly, it preserves all inner products. For an orthonormal basis,
\begin{equation}
\bra{j}U^\dagger(t)U(t)\ket{i}=\delta_{ji},
\end{equation}
and therefore
\begin{equation}
U^\dagger(t)U(t)=I.
\label{eq:l5-unitarity}
\end{equation}
The evolution operator is unitary.

At zero elapsed time \(U(0)=I\). Continuity then says that evolution through a small time \(\epsilon\) differs only slightly from the identity. In units with \(\hbar=1\), write
\begin{equation}
U(\epsilon)=I-i\epsilon H+O(\epsilon^2),
\qquad
U^\dagger(\epsilon)=I+i\epsilon H^\dagger+O(\epsilon^2).
\end{equation}
Substitution into Equation~\eqref{eq:l5-unitarity} gives
\begin{equation}
H^\dagger=H.
\end{equation}
The generator of time evolution is Hermitian and is therefore an observable. It is the Hamiltonian.

Applying the infinitesimal propagator to a state gives
\begin{equation}
\ket{\psi(\epsilon)}
=\ket{\psi(0)}-i\epsilon H\ket{\psi(0)}+O(\epsilon^2).
\end{equation}
Rearranging produces the difference quotient displayed on the board:
\begin{equation}
\frac{\ket{\psi(\epsilon)}-\ket{\psi(0)}}{\epsilon}
=-iH\ket{\psi(0)}+O(\epsilon).
\end{equation}

\begin{figure}[htbp]
\centering
\includegraphics[width=0.86\textwidth]{lecture_05_figure_03.png}
\caption{The infinitesimal update written as a difference quotient.}
\lectureframe{00:23:33}
\end{figure}

Taking \(\epsilon\to0\) yields the time-dependent Schr\"odinger equation:
\begin{equation}
\frac{d}{dt}\ket{\psi(t)}=-iH\ket{\psi(t)}.
\label{eq:l5-schrodinger}
\end{equation}
With constants restored, this is \(i\hbar\,d\ket{\psi}/dt=H\ket{\psi}\).

Schr\"odinger's original equation concerned the wavefunction of a moving particle. The operator statement above was once called the generalized Schr\"odinger equation because it applies equally to a spin, a three-state system, two spins, or a system of arbitrary complexity. The adjective has largely been dropped; this is now simply the general form of the equation.

Equation~\eqref{eq:l5-schrodinger} gives a definite evolution of the state when the system is left undisturbed. This does not make every experimental outcome definite. A quantum state determines probabilities for possible measurements, and the equation updates those probabilities through the changing amplitudes. Measurement generally changes the state and interrupts this isolated evolution.

\begin{classroomqa}
\audiencequestion{Does every measurement disturb the system?}
\lecturerresponse{Generically, yes, but there is an important special case. If the state is already an eigenstate of the observable, an ideal repeatable measurement can confirm its value without changing that state. What fails in general is the possibility of measuring an arbitrary observable without disturbing other, incompatible properties.}
\end{classroomqa}

\section{Energy and the Time-Independent Equation}

Because \(H\) is Hermitian, its eigenvalues are real and its eigenvectors can be chosen as a complete orthonormal basis:
\begin{equation}
H\ket{E_i}=E_i\ket{E_i}.
\label{eq:l5-time-independent}
\end{equation}
This is the time-independent form of the Schr\"odinger equation. It is an eigenvalue problem rather than an equation of motion: it asks which energy values the system can have and which states have those values sharply.

The spectrum depends on the system. Bound atomic motion and the harmonic oscillator have countable energy levels. An unbound particle may have a continuous spectrum. Energy can be measured in many ways; in principle even a sufficiently precise mass measurement could use \(E=mc^2\), although that is rarely the practical method. A real system already has its Hamiltonian; the physicist does not dictate it but tries to discover it. Once that operator is known, finding its eigenvalues and eigenvectors is the mathematical task.

Adding a constant \(cI\) to \(H\) shifts every energy by \(c\). It multiplies the whole state by one common phase \(e^{-ict}\), which does not change probabilities or expectation values. Only energy differences affect the dynamics considered here.

\section{How Expectation Values Move}

The standard phrase \emph{expectation value} means a statistical average. For an observable \(L\),
\begin{equation}
\expect{L}(t)=\bra{\psi(t)}L\ket{\psi(t)}.
\end{equation}
Only the state vectors carry time dependence at this stage. Differentiating both sides and using Equation~\eqref{eq:l5-schrodinger} together with its Hermitian conjugate,
\begin{equation}
\frac{d}{dt}\ket{\psi}=-iH\ket{\psi},
\qquad
\frac{d}{dt}\bra{\psi}=i\bra{\psi}H,
\end{equation}
gives
\begin{align}
\frac{d}{dt}\expect{L}
&=i\bra{\psi}HL\ket{\psi}
  -i\bra{\psi}LH\ket{\psi}\\
&=-i\expect{[L,H]}.
\label{eq:l5-expectation-motion}
\end{align}

\begin{figure}[htbp]
\centering
\includegraphics[width=0.86\textwidth]{lecture_05_figure_04.png}
\caption{The expectation-value equation of motion in commutator form.}
\lectureframe{00:35:35}
\end{figure}

The factor \(i\) is not a threat to the reality of \(\expect{L}\). For Hermitian \(L\) and \(H\), the commutator is anti-Hermitian:
\begin{equation}
[L,H]^\dagger=-[L,H].
\end{equation}
Its expectation value is therefore purely imaginary, so multiplying by \(-i\) gives a real derivative.

One sometimes writes the shorthand
\begin{equation}
\dot L=-i[L,H].
\end{equation}
At this point it should be read as an abbreviation for Equation~\eqref{eq:l5-expectation-motion}, an equation about averages. A full change of picture in which operators themselves carry the time dependence will come later.

\subsection{The Algebra of Commutators}

Commutators are antisymmetric:
\begin{equation}
[A,B]=-[B,A].
\end{equation}
They also obey a product rule analogous to the Leibniz rule for derivatives:
\begin{equation}
[AB,C]=A[B,C]+[A,C]B.
\label{eq:l5-product-rule}
\end{equation}
To verify it, preserve the order of every factor:
\begin{align}
A[B,C]+[A,C]B
&=A(BC-CB)\\
&\quad +(AC-CA)B\\
&=ABC-ACB\\
&\quad +ACB-CAB\\
&=ABC-CAB\\
&=[AB,C].
\end{align}

\begin{figure}[htbp]
\centering
\includegraphics[width=0.86\textwidth]{lecture_05_figure_05.png}
\caption{The commutator product rule checked by expanding both sides.}
\lectureframe{00:43:10}
\end{figure}

\begin{classroomqa}
\audiencequestion{Does the right side of the product rule accidentally count the left side twice?}
\lecturerresponse{No. On expansion, \(ABC\) and \(-CAB\) are the two desired terms. The two inserted terms are \(-ACB\) and \(+ACB\), which cancel exactly. The apparent extra copy is the temporary term needed to factor the expression in the two required ways.}
\end{classroomqa}

The same antisymmetry and product rule hold for classical Poisson brackets. This structural resemblance is not accidental. The classical correspondence is
\begin{equation}
\{A,B\}
\quad\longleftrightarrow\quad
\frac{1}{i\hbar}[A,B]
=-\frac{i}{\hbar}[A,B].
\label{eq:l5-poisson-correspondence}
\end{equation}
In \(\hbar=1\) units, the classical equation
\begin{equation}
\frac{dL}{dt}=\{L,H\}
\end{equation}
has the same form as Equation~\eqref{eq:l5-expectation-motion}. In ordinary macroscopic units \(\hbar\) is tiny, and commutators become negligible in the appropriate classical limit. This is the essential check that identifies the quantum generator \(H\) with the Hamiltonian of classical mechanics.

The factor of \(\hbar\) is required dimensionally. A Poisson bracket differentiates with respect to conjugate coordinates and momenta, whereas an operator commutator is built directly from products. Planck's constant supplies the conversion between those dimensions.

Set \(L=H\). Since every operator commutes with itself,
\begin{equation}
\frac{d}{dt}\expect{H}
=-i\expect{[H,H]}=0.
\end{equation}
For the time-independent Hamiltonians under discussion, the average energy is conserved. A stronger statement about the full energy distribution is possible, but the average is all that is needed here.

\begin{classroomqa}
\audiencequestion{What is \(L\) in these equations? Is it angular momentum?}
\lecturerresponse{Here \(L\) is a generic Hermitian linear operator, hence any observable. It may be a spin component, a position, or even \(H\) itself. The letter was chosen for \emph{linear operator} and should not be read as angular momentum unless that meaning is stated separately.}
\end{classroomqa}

\begin{classroomqa}
\audiencequestion{Why does the factor \(i\) not make a real average develop an imaginary derivative?}
\lecturerresponse{The commutator of two Hermitian operators is anti-Hermitian, so its expectation value is imaginary. The Pauli-matrix calculation below makes this concrete: their commutators contain an explicit factor \(2i\), which cancels the \(-i\) in the equation of motion.}
\end{classroomqa}

\section{Solving the Schr\"odinger Equation in the Energy Basis}

Assume the energy eigenvalue problem has been solved. Expand an arbitrary state in the fixed energy basis:
\begin{equation}
\ket{\psi(t)}=\sum_j\alpha_j(t)\ket{E_j}.
\label{eq:l5-energy-expansion}
\end{equation}
The basis vectors do not change in this description; all time dependence lies in the coefficients. Substituting Equation~\eqref{eq:l5-energy-expansion} into Equation~\eqref{eq:l5-schrodinger} gives
\begin{align}
\sum_j\dot\alpha_j(t)\ket{E_j}
&=-i\sum_j\alpha_j(t)H\ket{E_j}\\
&=-i\sum_j E_j\alpha_j(t)\ket{E_j}.
\end{align}
Uniqueness of the expansion in a basis implies, for each fixed \(j\),
\begin{equation}
\dot\alpha_j(t)=-iE_j\alpha_j(t),
\qquad\text{no sum on \(j\)}.
\end{equation}
The solution is
\begin{equation}
\alpha_j(t)=\alpha_j(0)e^{-iE_jt}.
\label{eq:l5-energy-phase}
\end{equation}

\begin{figure}[htbp]
\centering
\includegraphics[width=0.86\textwidth]{lecture_05_figure_06.png}
\caption{Each energy-basis coefficient satisfies its own exponential phase equation.}
\lectureframe{01:03:25}
\end{figure}

Thus the general solution is
\begin{equation}
\ket{\psi(t)}
=\sum_j\alpha_j(0)e^{-iE_jt}\ket{E_j}.
\label{eq:l5-general-solution}
\end{equation}
Restoring units replaces \(E_jt\) by \(E_jt/\hbar\). The magnitude of every coefficient stays fixed:
\begin{equation}
|\alpha_j(t)|=|\alpha_j(0)|.
\end{equation}
Time evolution changes relative phases. A single energy eigenstate acquires only one common phase and therefore has time-independent measurement probabilities.

The expectation value of \(L\) can now be written explicitly. Define its matrix elements in the energy basis by
\begin{equation}
L_{kj}=\bra{E_k}L\ket{E_j}.
\end{equation}
Then
\begin{align}
\expect{L}(t)
&=\sum_{k,j}
\alpha_k^*(0)\alpha_j(0)
e^{i(E_k-E_j)t}
L_{kj}.
\label{eq:l5-heisenberg-matrix}
\end{align}

\begin{figure}[htbp]
\centering
\includegraphics[width=0.86\textwidth]{lecture_05_figure_07.png}
\caption{The energy-basis matrix formula: every entry carries a transition-frequency phase.}
\lectureframe{01:10:45}
\end{figure}

The diagonal terms \(k=j\) are constant. An off-diagonal matrix element oscillates with the energy difference \(E_k-E_j\). In the historical account given here, Heisenberg recognized this matrix organization before the modern state-vector derivation was available; Dirac's notation later made its meaning transparent.

\begin{classroomqa}
\audiencequestion{Do the magnitudes of the \(\alpha_j\) coefficients remain fixed?}
\lecturerresponse{Yes. Each coefficient is multiplied by a complex phase of unit magnitude. Only relative phases change, and those phases are what allow expectation values containing off-diagonal matrix elements to vary with time.}
\end{classroomqa}

\begin{classroomqa}
\audiencequestion{Should the Schr\"odinger equation use a partial derivative or a total derivative?}
\lecturerresponse{For the abstract state \(\ket{\psi(t)}\) used here, time is the only displayed variable, so either notation carries the same content. Later a wavefunction such as \(\psi(x,t)\) will depend on both position and time, and the partial derivative notation will then be natural.}
\end{classroomqa}

\begin{classroomqa}
\audiencequestion{Is preservation of orthogonality merely an empirical assumption?}
\lecturerresponse{It is supported by experiment, but it is also unusually rigid. Attempts to deform unitary quantum mechanics tend to destroy positivity or conservation of probability, causal locality, or the consistent treatment of composite systems. David Finkelstein described deformable theories as undergoing plastic flow and rigid ones as brittle cracking. Quantum mechanics is brittle in this sense: small arbitrary changes do not usually leave a nearby consistent theory.}
\end{classroomqa}

\begin{classroomqa}
\audiencequestion{Why are the energy eigenvalues \(E_j\) real?}
\lecturerresponse{Unitarity implies that the infinitesimal generator \(H\) is Hermitian, and every eigenvalue of a Hermitian operator is real. The chain is therefore orthogonality preservation, unitarity, Hermiticity of \(H\), and reality of the energy spectrum.}
\end{classroomqa}

\section{A Spin in a Magnetic Field}

What Hamiltonian should be assigned to one spin? An isolated spin in otherwise empty space has no nontrivial preferred direction, so its Hamiltonian is proportional to the identity and produces no observable motion. A magnetic field supplies both an interaction energy and an axis.

Take the field along \(z\). A charged particle with spin carries a magnetic moment, and the familiar dipole energy is proportional to the alignment of that moment with the field. Absorb the field strength, magnetic moment, sign convention, and \(\hbar\) into one constant \(\omega\):
\begin{equation}
H=\frac{\omega}{2}\sigma_z.
\label{eq:l5-z-hamiltonian}
\end{equation}
The factor of \(1/2\) is convenient because Pauli commutators contain a factor of \(2\).

The classical arrow picture must be handled with care. A quantum spin does not possess simultaneously definite \(x\), \(y\), and \(z\) components, so no ordinary arrow can represent all of its properties at once. The useful three-dimensional object will be the vector of expectation values, not the two-component state vector itself.

\subsection{Pauli Commutators}

Use the matrix representations
\begin{equation}
\begin{aligned}
\sigma_x&=
\begin{pmatrix}0&1\\1&0\end{pmatrix},\\[0.4em]
\sigma_y&=
\begin{pmatrix}0&-i\\i&0\end{pmatrix},\\[0.4em]
\sigma_z&=
\begin{pmatrix}1&0\\0&-1\end{pmatrix}.
\end{aligned}
\end{equation}
For example,
\begin{align}
[\sigma_z,\sigma_x]
&=
\begin{pmatrix}1&0\\0&-1\end{pmatrix}
\begin{pmatrix}0&1\\1&0\end{pmatrix}
-
\begin{pmatrix}0&1\\1&0\end{pmatrix}
\begin{pmatrix}1&0\\0&-1\end{pmatrix}\\
&=
\begin{pmatrix}0&1\\-1&0\end{pmatrix}
-
\begin{pmatrix}0&-1\\1&0\end{pmatrix}\\
&=2
\begin{pmatrix}0&1\\-1&0\end{pmatrix}
=2i\sigma_y.
\end{align}
The cyclic relations are
\begin{equation}
\begin{aligned}
[\sigma_z,\sigma_x]&=2i\sigma_y,\\
[\sigma_x,\sigma_y]&=2i\sigma_z,\\
[\sigma_y,\sigma_z]&=2i\sigma_x.
\end{aligned}
\label{eq:l5-pauli-commutators}
\end{equation}
Reversing either entry reverses the sign.

\begin{figure}[htbp]
\centering
\includegraphics[width=0.86\textwidth]{lecture_05_figure_08.png}
\caption{The three Pauli commutators close cyclically on the Pauli matrices themselves.}
\lectureframe{01:33:30}
\end{figure}

These relations have the same cyclic structure as the classical Poisson brackets of angular momentum. They also close on the original three operators, which makes the equations of motion especially simple.

\begin{classroomqa}
\audiencequestion{Must the commutators be calculated with bra-ket brackets?}
\lecturerresponse{No. Operator multiplication is represented by matrix multiplication, so the \(2\times2\) matrices provide a convenient direct calculation. One could instead apply the operators successively to up and down basis states and obtain the same result.}
\end{classroomqa}

\subsection{Precession of the Expectation Values}

Apply Equation~\eqref{eq:l5-expectation-motion} with the Hamiltonian in Equation~\eqref{eq:l5-z-hamiltonian}:
\begin{align}
\frac{d}{dt}\expect{\sigma_x}
&=-i\frac{\omega}{2}\expect{[\sigma_x,\sigma_z]}
=-\omega\expect{\sigma_y},\\
\frac{d}{dt}\expect{\sigma_y}
&=-i\frac{\omega}{2}\expect{[\sigma_y,\sigma_z]}
=+\omega\expect{\sigma_x},\\
\frac{d}{dt}\expect{\sigma_z}
&=-i\frac{\omega}{2}\expect{[\sigma_z,\sigma_z]}
=0.
\label{eq:l5-spin-motion}
\end{align}

\begin{figure}[htbp]
\centering
\begin{minipage}[c]{0.62\textwidth}
\centering
\includegraphics[width=\textwidth]{lecture_05_figure_09.png}
\end{minipage}\hfill
\begin{minipage}[c]{0.31\textwidth}
\centering
\begin{tikzpicture}[scale=0.78,>=Latex]
  \draw[->,thick] (0,0) -- (0,3.0) node[above] {\(z\parallel\bm B\)};
  \draw[dashed] (0,1.8) ellipse (1.45 and 0.42);
  \draw[->,very thick] (0,0) -- (1.05,1.8)
    node[right] {\(\expect{\bm\sigma}\)};
  \draw[->,thick] (1.38,1.67)
    arc[start angle=-12,end angle=245,x radius=1.45,y radius=0.42];
\end{tikzpicture}
\end{minipage}
\caption{The board equations and their geometric consequence: the transverse averages rotate while the field-parallel component remains fixed. The arrow depicts the three expectation values, not a hidden classical spin.}
\lectureframe{01:35:28}
\end{figure}

The first two equations describe circular motion in the transverse plane:
\begin{align}
\expect{\sigma_x}(t)
&=\expect{\sigma_x}(0)\cos\omega t
-\expect{\sigma_y}(0)\sin\omega t,\\
\expect{\sigma_y}(t)
&=\expect{\sigma_x}(0)\sin\omega t
+\expect{\sigma_y}(0)\cos\omega t.
\end{align}
Meanwhile \(\expect{\sigma_z}\) is constant. The expectation-value vector therefore precesses around the magnetic-field axis with angular frequency \(\omega\).

The operators in this Schr\"odinger-picture calculation are fixed matrices. It is the state that changes, and hence the averages that execute the familiar gyroscopic motion.

\begin{classroomqa}
\audiencequestion{How can a two-dimensional state space produce a three-dimensional precession picture?}
\lecturerresponse{The state vector has two complex components, up to normalization and overall phase. The plotted real-space vector is instead \(\bigl(\expect{\sigma_x},\expect{\sigma_y},\expect{\sigma_z}\bigr)\). These expectation values are determined by the state and form the Bloch-vector picture, but they are not the state vector itself.}
\end{classroomqa}

\section{The Most General One-Spin Hamiltonian}

Every Hermitian \(2\times2\) matrix can be expanded as
\begin{equation}
H=aI+b_x\sigma_x+b_y\sigma_y+b_z\sigma_z,
\end{equation}
with real coefficients. The identity term changes only the overall phase. The three remaining coefficients can be regarded as the components of a vector. Writing their magnitude as \(\omega/2\) and their direction as a unit vector \(\bm n\), the nontrivial Hamiltonian becomes
\begin{equation}
H=\frac{\omega}{2}\,\bm\sigma\cdot\bm n,
\qquad
\bm\sigma\cdot\bm n
=n_x\sigma_x+n_y\sigma_y+n_z\sigma_z.
\label{eq:l5-general-spin-h}
\end{equation}
If \(\bm B=B\bm n\), then \(\bm n\) carries the direction of the field while \(\omega\) carries its magnitude together with the magnetic coupling. It is \(\bm n\), not \(\omega\), that records orientation.

The equation of motion takes the compact form
\begin{equation}
\frac{d}{dt}\expect{\bm\sigma}
=\omega\,\bm n\times\expect{\bm\sigma}.
\end{equation}
The component along \(\bm n\) commutes with \(H\) and remains fixed; the perpendicular components precess around \(\bm n\).

\begin{classroomqa}
\audiencequestion{Why is this the correct form of the Hamiltonian?}
\lecturerresponse{There are only four independent Hermitian \(2\times2\) matrices: \(I,\sigma_x,\sigma_y,\sigma_z\). The identity contributes no observable one-spin motion, so every nontrivial time-independent one-spin Hamiltonian is a real linear combination of the three Pauli matrices. A magnetic field packages those three coefficients into a magnitude and a direction.}
\end{classroomqa}

\begin{classroomqa}
\audiencequestion{What exactly does \(\omega\) represent?}
\lecturerresponse{It is a scalar proportional to the field magnitude times the magnetic moment, with conventional constants absorbed into it. It fixes both the precession rate and the energy splitting. The unit vector \(\bm n\) specifies the field direction.}
\end{classroomqa}

\subsection{The Two Energy Levels}

Writing Equation~\eqref{eq:l5-general-spin-h} as a matrix gives
\begin{equation}
\bm\sigma\cdot\bm n
=
\begin{pmatrix}
n_z & n_x-in_y\\
n_x+in_y & -n_z
\end{pmatrix}.
\end{equation}
Let its eigenvalues be \(\lambda_1,\lambda_2\). Its trace is zero:
\begin{equation}
\lambda_1+\lambda_2=0.
\end{equation}
For a unit vector \(n_x^2+n_y^2+n_z^2=1\), its determinant is
\begin{align}
\lambda_1\lambda_2
&=\det(\bm\sigma\cdot\bm n)\\
&=-n_z^2-(n_x-in_y)(n_x+in_y)\\
&=-(n_x^2+n_y^2+n_z^2)\\
&=-1.
\end{align}
The two eigenvalues are therefore
\begin{equation}
\lambda_\pm=\pm1.
\end{equation}

\begin{figure}[htbp]
\centering
\includegraphics[width=0.86\textwidth]{lecture_05_figure_10.png}
\caption{Trace and determinant fix the eigenvalues of \(\bm\sigma\cdot\bm n\) to \(+1\) and \(-1\).}
\lectureframe{01:59:20}
\end{figure}

Consequently,
\begin{equation}
E_\pm=\pm\frac{\omega}{2},
\qquad
\Delta E=E_+-E_-=\omega.
\label{eq:l5-spin-gap}
\end{equation}
The direction of the field changes the energy eigenvectors but not these two eigenvalues.

\begin{classroomqa}
\audiencequestion{How is photon emission related to these levels?}
\lecturerresponse{With the radiation field included, the higher spin level can decay to the lower level by emitting a photon of energy \(\Delta E=\omega\) in the present \(\hbar=1\) convention. The spin then flips between the two opposite orientations along the field axis. The precession calculation omitted this weak radiative coupling, so a spin may precess many times before emitting.}
\end{classroomqa}

\begin{classroomqa}
\audiencequestion{Why is the magnetic field treated as a number rather than a quantum operator?}
\lecturerresponse{The magnet is provisionally treated as a heavy classical apparatus whose fluctuations and recoil are negligible. A complete description would include it in a larger quantum system. The practical system-apparatus split is justified when the apparatus has many degrees of freedom and its relevant variables fluctuate so little that the electron cannot appreciably change them.}
\end{classroomqa}

\begin{classroomqa}
\audiencequestion{Do the smooth precession equations describe each individual trial?}
\lecturerresponse{They describe expectation values, hence averages over identically prepared trials. Individual measurements remain probabilistic. Repeated samples tend to approximate the theoretical average, but probability theory does not guarantee one exact result for every finite sequence, just as a million coin tosses need not contain exactly half heads.}
\end{classroomqa}

\begin{classroomqa}
\audiencequestion{If the spin points sideways, does a transition emit half a photon?}
\lecturerresponse{No. A sideways spin state is a superposition of the two energy eigenstates. In the symmetric case it has one-half probability in the upper level and one-half in the lower. It may therefore emit one full photon of energy \(\omega\), or emit none. Its average emitted energy can be \(\omega/2\), but no individual photon carries half the level spacing.}
\end{classroomqa}

We have now exhausted the nontrivial time-independent dynamics of a single spin. The same algebra that first signaled uncertainty also generates its motion: noncommuting Pauli matrices turn the Hamiltonian into precession, while the component commuting with the Hamiltonian remains conserved. The next increase in complexity is not another one-spin Hamiltonian but a larger state space, where independent systems can be composed and entangled.
